% Source: https://artofproblemsolving.com/community/c3h2122939_converging_vs_nonconverging_1__2__3___series_debate
% Stable identity: aops-topic-2122939-q01
\textbf{(a)} $S_n=\dfrac{n(n+1)}2\ge n$, so $S_n>M$ whenever $n>M$. The partial sums grow without bound, so the series diverges.

\textbf{(b)} The argument treats $S$ as a real number and then subtracts $9S$ from both sides. By (a) there is no such real number, so the equation $S=1+9S$ is meaningless and the rules of algebra do not apply to it. (Writing $S=\infty$ does not help either: $\infty$ is not a real number, and ``$\infty-\infty$'' is undefined.) What the grouping really shows is a true statement about partial sums: since $(3m-1)+3m+(3m+1)=9m$,
\[S_{3m+1}=1+9S_m\qquad\text{for every }m\ge1,\]
and both sides tend to infinity together; nothing like $-\dfrac18$ follows.

\textbf{(c)} Treating $T$ as a number: $T=1+(1+1+\cdots)=1+T$ gives $0=1$; and $T=(1+1)+(1+1)+\cdots=2T$ gives $T=0$, contradicting $T=1+T$. So manipulating divergent series like numbers can ``prove'' anything. (The value $-\dfrac1{12}$ sometimes attached to $1+2+3+\cdots$ comes from a different procedure, zeta-function regularisation, and is not a sum in the usual sense.) $\blacksquare$
